Adding one grain to a sandpile can leave it stable or trigger an avalanche across many sites. This project uses the two-dimensional Bak--Tang--Wiesenfeld model to investigate how that variation arises from a fixed local rule. The completed study compares uniform input with spatially biased input, cylindrical boundaries and four-grain batches. A queue simulator checked against independent implementations processed \R{productionGrains} production grains across four baseline sizes and five seeds per condition. At $L=64$, the uniform open model had mean height \R{base64height} and mean topplings \R{base64S_grain} per grain, compared with the finite-grid prediction \R{base64prediction}. Bias and cylindrical boundaries increased mean topplings to \R{bias64S_grain} and \R{cyl64S_grain}, while mean loss remained close to one per grain. A small height difference under biased input was not reproduced in separate, longer runs. All \R{matchedBatches} matched batches reproduced grouped sequential topplings and loss exactly, although their event distributions changed. Normalized discrete fits favour cutoff descriptions over pure powers on the chosen finite ranges. These comparisons distinguish changes in the stationary state measure, the response to input and the grouping of avalanche events.
